% ======================================================================
% 第6章 評価実験
% ======================================================================

% 【追加】章の導入（位置付け）
本章では，前章で提案したRestart PSOの有効性を検証するために実施したシミュレーション実験の結果について述べる．
比較対象として厳密解法および従来手法（Global PSO等）を用い，第1章で述べた計算速度，解の精度，および厳しい制約条件下でのロバスト性という3つの観点から，提案手法の優位性を定量的かつ多角的に評価する．

\section{概要}
実験は，PythonおよびNetworkXを用いて実装したシミュレータ上で行い，以下の3つの評価軸に基づいて提案手法の特性を検証した．
各評価軸に対応する具体的な実験設定を表 \ref{tab:exp_summary} に示す．

\begin{table}[tb]
  \centering
  \caption{評価実験の構成と対応する節}
  \label{tab:exp_summary}
  \begin{tabular}{l|l|l} \hline
    実験ID & 実験項目 & 評価結果の記載箇所 \\ \hline
    実験1 & 並列化スケーラビリティ検証 & \ref{sec:exp_speed}節 (図\ref{fig:parallel_time}) \\
    実験2 & ノード数・制約数と計算時間の関係 & \ref{sec:exp_speed}節 (図\ref{fig:time_vs_nodes_1const}, \ref{fig:time_vs_nodes_3const}) \\
    実験3 & 厳密解に対する近似率（解精度） & \ref{sec:exp_accuracy}節 (図\ref{fig:accuracy_vs_nodes_1const} $\sim$ \ref{fig:accuracy_comparison}) \\
    実験4 & 世代ごとの収束特性 & \ref{sec:exp_accuracy}節 (図\ref{fig:convergence}) \\
    実験5 & 厳しい制約下でのロバスト性 & \ref{sec:exp_robustness}節 (図\ref{fig:feasibility_rate}, \ref{fig:pareto_count}) \\ \hline
    実験6 & エンコーディング処理のボトルネック解析 & \ref{sec:exp_bottleneck}節 (図\ref{fig:encoding_linear}, \ref{fig:encoding_log}) \\ \hline
  \end{tabular}
\end{table}

\section{実験環境とパラメータ設定}
シミュレーション実験は，表 \ref{tab:env} に示す計算機環境にて実施した．

\begin{table}[tb]
  \centering
  \caption{シミュレーション環境}
  \label{tab:env}
  \begin{tabular}{l|l} \hline
    OS & macOS \\
    CPU & Apple M2 (8 コア) \\
    メモリ & 8 GB \\
    使用言語 & Python 3.12 \\
    ライブラリ & NetworkX, NumPy, Multiprocessing \\ \hline
  \end{tabular}
\end{table}

本実験では，ネットワークトポロジとして，ノード間の接続確率を0.2としたランダムグラフを採用した．
各リンク（エッジ）には，ボトルネック帯域（Bandwidth），遅延（Delay），パケット損失率（Loss Rate），および信頼性（Reliability）の4つのQoS属性を付与した．
シミュレーションで使用したトポロジ生成パラメータおよび各属性の設定範囲を表 \ref{tab:topology_params} に示す．

\begin{table}[tb]
  \centering
  \caption{ネットワークトポロジとQoS属性の設定値}
  \label{tab:topology_params}
  \begin{tabular}{l|l} \hline
    項目 & 設定内容 / 範囲 \\ \hline
    トポロジモデル & ランダムグラフ \\
    接続確率 & 0.2 \\
    リンク帯域幅 (Bandwidth) & $1 \sim 50$ (一様乱数) \\
    リンク遅延 (Delay) & $1.0 \sim 10.0$ ms (一様乱数) \\
    パケット損失率 (Loss) & $0.01\% \sim 3.0\%$ (一様乱数) \\
    信頼性 (Reliability) & $99.0\% \sim 99.999\%$ (一様乱数) \\ \hline
  \end{tabular}
\end{table}

また，多制約経路問題（MCOP）における制約条件 $C$ は，アプリケーションの要求品質を模して以下のように設定した．
遅延制約 $D_{\max}$ については，ネットワーク形状により変動する最短遅延 $D_{\min}$ の3.0倍を許容上限とした ($D_{\max} = 3.0 \times D_{\min}$)．
パケット損失率制約 $L_{\max}$ は $10\%$ ($0.1$) 以下，信頼性制約 $R_{\min}$ は $95\%$ ($0.95$) 以上とした．
なお，損失率と信頼性については，計算上対数変換を用いた加法的なコストとして処理している．

本研究におけるPSOのパラメータ設定について述べる．
すべての実験において共通して使用した係数パラメータを表 \ref{tab:common_params} に示す．
慣性重み $w$ および加速係数 $c_1, c_2$ は，探索初期の大域的探索と終盤の局所探索のバランスを考慮し，世代数に応じて線形に変化させる時変パラメータを採用した．
この動的パラメータ戦略（Adaptive）の選定根拠となった予備実験の結果については，付録\ref{app:parameter_tuning}に示す．
\begin{table}[tb]
  \centering
  \caption{PSOの共通パラメータ}
  \label{tab:common_params}
  \begin{tabular}{l|l} \hline
    パラメータ & 設定値 \\ \hline
    慣性重み ($w$) & $0.9 \to 0.18$ (線形減少) \\
    認知係数 ($c_1$) & $3.32 \to 0.44$ \\
    社会係数 ($c_2$) & $0.88 \to 3.66$ \\
    リスタート閾値 & 20世代停滞 \\ \hline
  \end{tabular}
\end{table}

また，実験の目的に応じて設定した粒子数，最大世代数，および探索の終了条件を表 \ref{tab:exp_settings} にまとめる．
特に，大規模ネットワークにおける計算時間の計測である実験2においては，大域解 $gBest$ が50世代連続で更新されなかった時点で収束とみなし，探索を早期終了する設定とした．
その他の実験1，3，4および5については，収束過程や成功率を詳細に分析するため，最大世代数まで探索を実行した．

\begin{table}[tb]
  \centering
  \caption{実験ごとの個別設定}
  \label{tab:exp_settings}
  \begin{tabular}{l|c|c|l} \hline
    実験項目 & 粒子数 & 最大世代数 & 終了条件 \\ \hline
    実験1 & 100 & 100 & 固定世代数 \\
    実験2 & 100 & 100 & 50世代停滞による収束判定 \\
    実験3 & 200 & 200 & 固定世代数 \\
    実験4 & 150 & 100 & 固定世代数 \\
    実験5 & 100 & 100 & 固定世代数 \\ \hline
  \end{tabular}
\end{table}

本実験では，提案手法の性能を多角的に評価するため，以下の3つの手法と比較を行った．

\begin{itemize}
  \item {Global PSO}: 
  従来の標準的なPSOであり，群れ全体で単一の大域最良解 ($gBest$) を共有する手法である．情報の伝播が速いが，複雑な多峰性関数においては局所解に陥りやすい特性を持つ．
  
  \item {Local PSO}:
  提案手法と同じく空間的近傍を用いた $lBest$ トポロジーを採用しているが，リスタート（Explosion）戦略を実装していない手法である．
  提案手法 (Restart PSO) はこれにリスタート戦略を加えたものであり，両者を比較することでリスタート機能単体による性能向上効果を分離して評価する．
  
  \item {厳密解法}: 
  ラベル修正法に基づき，パレート最適解を網羅的に探索する手法である．計算コストは高いが真の最適解を導出できるため，解の精度（近似率）を算出する際の基準値として用いる．
\end{itemize}

\section{並列化による計算時間短縮効果}
\label{sec:exp_speed}
本節では，表\ref{tab:exp_summary}における実験1および実験2の結果に基づき，提案手法の計算速度について評価する．

まず，プロセス並列化が計算時間の短縮に有効であることを示す．
図 \ref{fig:parallel_time} に，ノード数1000のネットワークにおいて並列処理のコア数を1から4まで変化させた際の平均計算時間の推移を示す．

\begin{figure}[tb]
  \centering
  \includegraphics[width=11cm]{images/parallel_time.png}
  \caption{並列化による平均計算時間の推移（$N=1000$）}
  \label{fig:parallel_time}
\end{figure}

図 \ref{fig:parallel_time} より，コア数の増加に伴い計算時間が顕著に短縮されていることが確認できる．
具体的には，4コア並列化を行うことで計算時間は1コア時の約1/3程度にまで短縮されている．
提案手法（Restart PSO）は，リスタート処理などの追加演算を含むため，1コア時の計算時間は単純なGlobal PSOよりも長い．
しかし，並列化によってオーバーヘッドを相殺し，4コア使用時には1コア動作時のGlobal PSOよりも短い時間で処理を完了できている．なお，並列化に伴うプロセス間通信（データ転送等）のオーバーヘッドは存在するが，本手法では各粒子における経路生成と評価の計算負荷が支配的であるため，オーバーヘッドの影響を上回る十分な速度向上が得られている．
これにより，提案手法の並列実装がリアルタイム性を確保する上で効果的に機能していると言える．

次に，ネットワーク規模および制約数の増加に対するスケーラビリティについて述べる．
ノード数を100から4000まで変化させた際の厳密解法との比較結果を，図 \ref{fig:time_vs_nodes_1const} （制約数1）および図 \ref{fig:time_vs_nodes_3const} （制約数3）に示す．

\begin{figure}[tb]
  \centering
  \includegraphics[width=11.0cm]{images/mini計算時間比較_1制約3手法.png}
  \caption{ノード数に対する平均計算時間の比較（制約数1）}
  \label{fig:time_vs_nodes_1const}
\end{figure}

\begin{figure}[tb]
  \centering
  \includegraphics[width=11.2cm]{images/mini計算時間比較_3制約3手法.png}
  \caption{ノード数に対する平均計算時間の比較（制約数3）}
  \label{fig:time_vs_nodes_3const}
\end{figure}

図 \ref{fig:time_vs_nodes_1const} および図 \ref{fig:time_vs_nodes_3const} の比較から，提案手法は厳密解法と比較して制約数の増加に対する計算時間の変動が極めて小さいことがわかる．
厳密解法は，制約数が1つの場合に比べて3つの場合では計算時間が指数関数的に増大しており，大規模ネットワークへの適用が困難であることを示している．
これは，制約数の増加に伴いパレート最適解の数が増え，探索すべきラベルの数が爆発的に増加するためである．
一方，提案手法を含むPSOは，制約数が増加しても計算時間の変化は微小であった．
これは，PSOの計算負荷が主に粒子数と世代数に依存し，制約条件の数自体には大きく依存しないアルゴリズム特性によるものである．
以上の結果から，特に複数のQoS制約を考慮する必要がある大規模ネットワーク環境下において，提案手法は厳密解法よりも優れたスケーラビリティを有していると結論付けられる．

\section{解の探索性能と精度の評価}
\label{sec:exp_accuracy}
本節では，実験3および実験4の結果に基づき，提案手法が発見する解の品質について評価する．

まず，ネットワーク規模が増大しても提案手法が高い解精度を維持できることを示す．
図 \ref{fig:accuracy_vs_nodes_1const} および図 \ref{fig:accuracy_vs_nodes_3const} に，ノード数を変化させた際の厳密解に対する近似率の推移を示す．

\begin{figure}[b]
  \centering
  \includegraphics[width=11cm]{images/mini1制約解精度比較_Global_vs_Restart.png}
  \caption{ノード数に対する解精度の推移（制約数1）}
  \label{fig:accuracy_vs_nodes_1const}
\end{figure}

\begin{figure}[tb]
  \centering
  \includegraphics[width=11cm]{images/mini3制約解精度比較_Global_vs_Restart.png}
  \caption{ノード数に対する解精度の推移（制約数3）}
  \label{fig:accuracy_vs_nodes_3const}
\end{figure}

図より，Global PSOはノード数が増加するにつれて近似率が低下する傾向が見られる．これは，探索空間が広がるにつれて局所解に陥りやすくなるためである．
対して，提案手法（Restart PSO）はノード数が増加しても一定の高い近似率を維持している．
これは，リスタート戦略によって探索の停滞を回避し，広大な探索空間の中からでもより良い解を発見できていることを示唆している．

さらに，標準的な制約条件（遅延倍率3.0，$N=1000$）において各手法が到達した最終的な解精度を比較した結果を図 \ref{fig:accuracy_comparison} に示す．

\begin{figure}[tb]
  \centering
  % 画像ファイルのパスを指定 (例: imagesフォルダの中にある場合)
  \includegraphics[width=10cm]{images/卒論用3手法解精度par200gen200.png}
  \caption{各手法における近似率の比較 ($N=1000$, 粒子数=200, 世代数=200)}
  \label{fig:accuracy_comparison}
\end{figure}

図 \ref{fig:accuracy_comparison} において，提案手法は近似率0.888を達成し，Global PSOの0.599やLocal PSOの0.822と比較して最も高い精度を示した．
Global PSOは早期収束により探索を打ち切ってしまう傾向があるが，提案手法はリスタート機能により探索を継続し，精度の向上を図れていることが数値的にも裏付けられた．

また，提案手法が世代経過とともに解を改善していく様子を図 \ref{fig:convergence} に示す．

\begin{figure}[tb]
  \centering
  \includegraphics[width=11.5cm]{images/par150gen100restart10gen.png}
  \caption{世代数に対するボトルネック帯域幅の改善推移}
  \label{fig:convergence}
\end{figure}

図 \ref{fig:convergence} より，リスタートが発生したタイミング（グラフ上の変曲点）以降も，解の改善が継続していることが確認できる．
これは，一度局所解に陥った後でも，リスタートによって群れ全体の探索ベクトルが再設定され，新たな有望領域を発見できていることを示している．
したがって，提案手法は単に速いだけでなく，十分な世代数を重ねることで厳密解に近い高品質な経路を導出可能であると言える．

\section{厳しい制約条件下でのロバスト性}
\label{sec:exp_robustness}
最後に，実験5の結果に基づき，探索空間が狭く解の発見が困難な状況における提案手法の有効性を検証する．

遅延制約の倍率を変化させながら，各手法が100回の試行中で何回制約を満たす解（実行可能解）を発見できたか（成功率）を計測した結果を，図 \ref{fig:feasibility_rate} に示す．
倍率が小さいほど，制約が厳しく実行可能解の数が少ないことを意味する．

\begin{figure}[tb]
  \centering
  % 画像ファイルのパスを指定 (例: imagesフォルダの中にある場合)
  \includegraphics[width=11.5cm]{images/卒論用feasibility_rate_result_fix.png}
  \caption{遅延制約の厳しさと成功率の関係}
  \label{fig:feasibility_rate}
\end{figure}

図 \ref{fig:feasibility_rate} より，制約が厳しい領域（1.1倍〜1.2倍）において，手法間で明確な性能差が見られた．
Global PSOやLocal PSOは成功率が著しく低下しているのに対し，Restart PSOは比較的高い成功率を維持している．
これは，従来手法が一度の探索で実行可能解を見つけられずに収束してしまうような状況でも，提案手法はリスタートによって探索をやり直し，粘り強く実行可能解を探し出せる能力が高いことを示している．

この結果をさらに裏付けるため，厳密解法によって求められたパレート最適経路の数と，制約倍率の関係を図 \ref{fig:pareto_count} に示す．

\begin{figure}[tb]
  \centering
  % 画像ファイルのパスを指定 (例: imagesフォルダの中にある場合)
  \includegraphics[width=11.5cm]{images/pareto_count_analysis_fix.png}
  \caption{パレート最適解数（探索空間の広さ）と成功率の関係}
  \label{fig:pareto_count}
\end{figure}

図 \ref{fig:pareto_count} の赤線が示すように，制約倍率が1.1倍付近ではパレート最適解の数が極端に少なくなり，探索空間内の当たりが極めて少ない状態となっている．
このような探索が困難な状況においても，緑色の破線で示される提案手法が一定の成功率を保っていることは特筆すべき点である．
以上の結果から，提案手法は厳しいQoS要件が課される環境下においても，高い確率で実行可能解を発見できるロバスト性を有していると結論付けられる．

\section{エンコーディング処理のボトルネック解析}
\label{sec:exp_bottleneck}
本章これまでの実験により，提案手法の並列化が計算時間の短縮に大きく寄与していることが確認された．
本節では，この並列化の有効性を裏付けるために，アルゴリズム内部で最も計算負荷が高い優先度ベースエンコーディング処理単体の詳細な性能解析を行う．
グラフ密度やノード規模が計算時間に与える影響を分析することで，並列化アプローチの妥当性を検証する．

まず，ノード数 $N$ を100から10000まで変化させた際のエンコーディング処理の実行時間を計測した結果について述べる．
比較対象として，平均次数を10に固定した疎グラフと，本研究で採用している接続確率0.2の密グラフを用いた．
図\ref{fig:encoding_linear}に示す線形スケールでの比較結果を見ると，疎グラフ条件ではノード数が増加しても計算時間は緩やかな上昇に留まっている．
対照的に，密グラフ条件では，ノード数の増加に伴い計算時間が急激に増大している．
特に $N=2000$ の時点ですでに数倍の性能差が生じており，探索空間の膨大化が計算負荷の主要因であることを示している．

\begin{figure}[tb]
  \centering
  \includegraphics[width=11.5cm]{images/pathencode_1.png}
  \caption{グラフ密度によるエンコーディング時間の推移（線形スケール）}
  \label{fig:encoding_linear}
\end{figure}

次に，計算量のオーダーとハードウェア的な限界を分析するため，同データを対数スケールでプロットした結果を図\ref{fig:encoding_log}に示す．
図中の点線は理論的な計算量オーダー $O(N^2)$ を示している．
$N \le 2000$ の領域では理論通り $O(N)$ から $O(N^2)$ の範囲で推移しているが，$N \ge 3000$ の領域においてグラフの傾きが急変し，理論値である $O(N^2)$ すらも逸脱して悪化する傾向が確認された．

\begin{figure}[tb]
  \centering
  \includegraphics[width=11.5cm]{images/pathencode_taisutaisu.pdf}
  \caption{エンコーディング時間の対数スケール評価}
  \label{fig:encoding_log}
\end{figure}

この理論値からの乖離要因を特定するため，プロファイリングによる解析を行った．
その結果，処理時間の99\%以上がループ処理とリスト操作に集中していることが判明した．
特に大規模な密グラフでは，経路生成時に参照する隣接ノードリストのサイズが巨大化し，頻繁なメモリ確保とガベージコレクションが発生する．
加えて，リスト順に不連続なメモリ領域へアクセスする際のデータ量がCPUキャッシュ容量を超過し，低速なメインメモリへのアクセス待ちが多発していると考えられる．

以上の分析から，本手法におけるボトルネックは，アルゴリズムの計算量 $O(N^2)$ に加え，大規模データ処理に伴うメモリアクセスの物理的遅延にあることが明らかとなった．
1回あたりの遅延はわずかであっても，PSO全体では数万回の反復が行われるため，この遅延が致命的となる．
このボトルネックは単一プロセスのコード最適化だけでは解消が困難な物理的制約に起因するものであり，したがって，第\ref{sec:exp_speed}節で示したプロセス並列化によって計算とメモリ負荷を分散させるというアプローチは，大規模ネットワークへの適用において合理的かつ不可欠な解決策であると結論付けられる．